% !TEX program = pdflatex
\documentclass[11pt]{article}
\usepackage[T1]{fontenc}
\usepackage[utf8]{inputenc}
\usepackage[margin=1in]{geometry}
\usepackage{amsmath,amssymb,amsthm,mathtools}
\usepackage{bm}
\usepackage{booktabs}
\usepackage{microtype}
\usepackage{enumitem}
\usepackage{hyperref}
\usepackage{graphicx}
\usepackage{xcolor}
\usepackage{listings}
\lstset{basicstyle=\ttfamily\small,breaklines=true,columns=fullflexible,frame=single}
\hypersetup{colorlinks=true,linkcolor=blue,citecolor=blue,urlcolor=blue}

\title{The Alpha Function in IFMD: Formalization, Evaluation, and\\ Safety-Certified Integration}
\author{Ryan O. Van Gelder}
\date{\today}

% Theorem environments
\newtheorem{definition}{Definition}
\newtheorem{theorem}{Theorem}
\newtheorem{proposition}{Proposition}
\newtheorem{lemma}{Lemma}
\theoremstyle{remark}
\newtheorem{remark}{Remark}

% Shorthand
\DeclareMathOperator{\diag}{diag}
\newcommand{\R}{\mathbb{R}}
\newcommand{\C}{\mathbb{C}}
\newcommand{\E}{\mathbb{E}}
\newcommand{\Var}{\mathrm{Var}}
\newcommand{\RePart}{\operatorname{Re}}

\begin{document}
\maketitle

\vspace{0.5em}

\begin{abstract}
We formalize a unifying \emph{alpha function} \(\alpha(x;\theta)\) and integrate it with the IFMD stack (ACE, PETC, Langlands Prism, PIRTM). The report provides: (i) a precise integral/series definition with machine-checkable domain guards; (ii) certified numerical evaluation using Gauss--Laguerre quadrature and adaptive series; (iii) recovery of classical special functions (Gamma, Beta, Zeta, Bessel, generalized hypergeometric); (iv) a diagnostics schema suitable for safety logging; (v) kernel ablation with closed-form references; and (vi) a projection-first ACE integration pattern with certificates (GapLB, SlopeUB). A minimal evaluator implementation is included in Appendix~\ref{app:code}.
\end{abstract}

\newpage
\tableofcontents

\section{Introduction}
\label{sec:intro}
The IFMD architecture enables safe, mathematically grounded AI control. Its core layers are: ACE (safety), PETC (structure), the Langlands Prism (integration), and PIRTM (projection/runtime monitoring). The \emph{alpha function} consolidates a wide range of classical functions into a single parameterized kernel, enabling (a) unified numerical guarantees, (b) structured feature extraction, and (c) streamlined verification.

\paragraph{Design principles.} (1) Separate inspiration (for example, quantum or number-theoretic motifs) from certified control dynamics; (2) prefer kernels with closed-form checks; (3) attach diagnostics and guards to every evaluation; (4) enforce projection-first actuation via ACE.

\section{Formal Definition}
\label{sec:formal}
\begin{definition}[Alpha Master]
Let \(x>0\) and parameters \(\theta\in\Theta\subset\R^m\). Define
\begin{equation}
\label{eq:alpha-master}
\alpha(x;\theta)\;=\;\int_{0}^{\infty} t^{\theta_0-1}\,e^{-xt}\,\mathcal{G}(t;\theta)\,dt\; +\; \sum_{k\in\mathcal{K}(\theta)} c_k(\theta)\,x^{\rho_k(\theta)}.
\end{equation}
The kernel family \(\mathcal{G}(\cdot;\theta)\) is \emph{lawful} if it satisfies the domain and convergence conditions below. The discrete term is optional and used for series-style constructions.
\end{definition}

\paragraph{Domain and convergence guards.} A call to \(\alpha(x;\theta)\) is permitted only if all guards pass:
\begin{itemize}[nosep]
  \item \textbf{Integral path:} \(x>0\) and \(\RePart(\theta_0)>0\) with an \(\mathcal{G}\) such that \(\int_0^{\infty}\!\big|t^{\theta_0-1} e^{-xt}\,\mathcal{G}(t;\theta)\big|\,dt<\infty\).
  \item \textbf{Series path:} Discrete parameters ensure absolute convergence; for example, for \(\zeta(s)=\sum_{n\ge1}n^{-s}\), require \(\RePart(s)>1\).
  \item \textbf{Diagnostics:} Each evaluation returns an estimated error and convergence flags (Section~\ref{sec:diagnostics}).
\end{itemize}

\section{Recovery of Classical Functions}
\label{sec:recovery}
We collect parameter slices and constructions that recover standard special functions from~\eqref{eq:alpha-master}.

\subsection{Gamma and Beta}
\paragraph{Gamma.} With \(x=1\), \(\theta_0=s\), and \(\mathcal{G}\equiv1\),
\begin{equation}
\alpha(1;\theta)=\int_0^{\infty}t^{s-1}e^{-t}\,dt=\Gamma(s),\qquad \RePart(s)>0.
\end{equation}
\paragraph{Beta.} Using \(B(a,b)=\Gamma(a)\Gamma(b)/\Gamma(a+b)\) and the Gamma slice twice yields \(B(a,b)\) for \(\RePart(a),\RePart(b)>0\). A direct Beta-integral recovery is also possible via a change of variables onto \([0,1]\).

\subsection{Zeta, Bessel, and Hypergeometric}
\paragraph{Riemann zeta (real slice).} With discrete-only part \(\sum_{n\ge1} n^{-s}\), we require \(\RePart(s)>1\). An explicit truncation remainder bound is \(R_N\le N^{1-s}/(s-1)\) for real \(s>1\).

\paragraph{Bessel $J_\nu$.} The series recovery is
\begin{equation}
J_\nu(z)=\sum_{m\ge0}\frac{(-1)^m}{m!\,\Gamma(m+\nu+1)}\Big(\frac{z}{2}\Big)^{2m+\nu},\qquad \nu>-1.
\end{equation}

\paragraph{Generalized hypergeometric.} For parameters \(\bm{a}=(a_1,\ldots,a_p)\) and \(\bm{b}=(b_1,\ldots,b_q)\),
\begin{equation}
{}_pF_q(\bm{a};\bm{b};z)=\sum_{k\ge0}\frac{(a_1)_k\cdots(a_p)_k}{(b_1)_k\cdots(b_q)_k}\,\frac{z^k}{k!},\qquad (a)_k=\frac{\Gamma(a+k)}{\Gamma(a)}.
\end{equation}
Convergence holds for all \(z\) if \(q\ge p\); if \(q=p-1\) then \(|z|<1\).

\section{Numerical Evaluator}
\label{sec:evaluator}
\subsection{Gauss--Laguerre Quadrature (Integral Path)}
Using the change of variables \(u=xt\) gives
\begin{equation}
\alpha(x;\theta)=x^{-\theta_0}\int_0^{\infty} u^{\theta_0-1}e^{-u}\,\mathcal{G}(u/x;\theta)\,du.
\end{equation}
We approximate \(\int_0^{\infty} f(u)\,e^{-u}\,du\) via Gauss--Laguerre with adaptive node doubling \(N\in\{16,32,64\}\), terminating when successive estimates differ by less than a tolerance \(\varepsilon\).

\subsection{Series Control}
For discrete terms, we use adaptive truncation with ratio tests or explicit tail bounds (for example, zeta).

\subsection{Returned Diagnostics}
\label{sec:diagnostics}
Each evaluation returns a structured dictionary
\begin{equation*}
\texttt{diagnostics}=\{\,\texttt{convergence\_flags},\ \texttt{estimated\_error},\ \texttt{computation\_path},\ \texttt{performance\_metrics}\,\},
\end{equation*}
where
\begin{align*}
\texttt{convergence\_flags}&=\{\texttt{integral:bool},\ \texttt{series:bool}\},\\
\texttt{computation\_path}&\in\{\texttt{integral},\ \texttt{series},\ \texttt{hybrid}\},\\
\texttt{performance\_metrics}&=\{\texttt{nodes\_used:int},\ \texttt{terms\_used:int},\ \texttt{compute\_time:float}\}.
\end{align*}

\section{Kernel Ablation}
\label{sec:ablation}
We compare kernels with closed-form references:
\begin{description}[leftmargin=1.8em,labelsep=0.5em]
  \item[\(\bm{G_1}\): baseline] \(\mathcal{G}(t)=1\) gives \(\alpha=\Gamma(s)\,x^{-s}\).
  \item[\(\bm{G_2}\): exp-shift] \(\mathcal{G}(t)=e^{a t}\) with \(a<x\) gives \(\alpha=\Gamma(s)\,(x-a)^{-s}\).
  \item[\(\bm{G_3}\): polynomial] \(\mathcal{G}(t)=(1+ct)^m\) gives \(\alpha=\sum_{k=0}^m\binom{m}{k}c^k\,\Gamma(s+k)\,x^{-(s+k)}\).
\end{description}
Each run logs absolute/relative error, nodes used, and compute time.

\section{ACE Integration Pattern}
\label{sec:ace}
\subsection{Feature Extraction}
Sample \(x\) on a geometric grid and form a PETC-typed feature vector
\begin{equation}
\bm{\phi}=\mathrm{normalize}\Big( \big[\,\alpha(x_i;\theta),\ \partial_{\log x}\log\alpha(x_i;\theta),\ \partial_{\log x}^2\log\alpha(x_i;\theta)\,\big]_{i=1}^M \Big).
\end{equation}

\subsection{Projection-First Actuation}
Let \(\tilde{\bm{w}}\) be the estimator output. We project onto a weighted \(\ell_1\) safety set \(\mathcal{S}=\{\bm{w}:\sum_p b_p|w_p|\le T\}\).
\begin{proposition}[Soft-thresholding solves the weighted \(\ell_1\) projection]
For any \(\tilde{\bm{w}}\in\R^d\) and weights \(b_p>0\), there exists \(\lambda\ge0\) such that the solution of
\(\min_{\bm{w}}\tfrac12\|\bm{w}-\tilde{\bm{w}}\|_2^2\) subject to \(\sum_p b_p|w_p|\le T\) is \(w_p=\mathrm{sign}(\tilde w_p)\max(|\tilde w_p|-\lambda b_p,0)\) with \(\lambda\) chosen by bisection so \(\sum_p b_p|w_p|=T\) when the constraint is active.
\end{proposition}
\begin{remark}
Certificates log a contraction margin (GapLB), a Lipschitz envelope (SlopeUB), and feasibility. On failure, the controller rolls back.
\end{remark}

\section{PETC Typing and Lawfulness}
\label{sec:petc}
Prime-indexed or arithmetic features are ledgered with \emph{prime signatures}. Typed tensor axes ensure contractive operators preserve signatures under reshapes and contractions. Lawfulness budgets bound how much such features can influence proposals before ACE projection.

\section{Experiments and Results}
\label{sec:experiments}
\paragraph{E1: Recovery tests.} Compare numerically evaluated \(\alpha\) slices to references for Gamma, Beta, Zeta, Bessel, and \({}_pF_q\). Report max relative error and diagnostics.

\paragraph{E2: PETC feature check.} Feed \(\bm{\phi}\) into a small estimator; verify ACE projection maintains contraction on random plants. Record GapLB and SlopeUB.

\paragraph{E3: Kernel ablation.} Run \(\bm{G_1}\)--\(\bm{G_3}\) with references; log errors, nodes, and timings.

\section{Reproducibility}
\label{sec:reprod}
\begin{itemize}[nosep]
  \item \textbf{Python:} 3.11; packages: \texttt{numpy}, \texttt{pandas}.
  \item \textbf{Artifacts:} evaluator module, recovery CSVs, kernel ablation CSV.
  \item \textbf{Invocation:} code is included below as an excerpt for portability.
\end{itemize}

\appendix
\section{Notation}
\label{app:notation}
\(\Gamma(\cdot)\) Gamma function; \(B(\cdot,\cdot)\) Beta; \(\zeta(\cdot)\) Riemann zeta; \({}_pF_q\) generalized hypergeometric; \(J_\nu\) Bessel of the first kind.

\section{Minimal Evaluator Code}
\label{app:code}
\begin{lstlisting}[language=Python,caption=alpha_eval.py (excerpt)]
import math
import numpy as np
from numpy.polynomial.laguerre import laggauss

def gauss_laguerre_integrate(f, nodes=32):
    x, w = laggauss(nodes)
    return float(np.dot(w, f(x)))

def alpha_integral(x, s, G=lambda t: 1.0, tol=1e-10, max_nodes=64, min_nodes=16):
    assert x > 0 and s > 0
    def integrand(u):
        return (u**(s-1.0)) * G(u/x)
    nodes, prev = min_nodes, None
    x_pow = x**(-s)
    while nodes <= max_nodes:
        val = gauss_laguerre_integrate(integrand, nodes=nodes) * x_pow
        if prev is not None and abs(val - prev) < tol:
            return val
        prev, nodes = val, 2*nodes
    return prev
\end{lstlisting}

\section{Bibliography (internal working notes)}
\begin{thebibliography}{99}
\bibitem{Alpha} Alpha Function (working draft). Internal PDF.
\bibitem{ACE} Arithmetic Control Engine (ACE). Internal PDF.
\bibitem{PETC} Prime-Encoded Tensor Calculus (PETC). Internal PDF.
\bibitem{LanglandsPrism} Langlands Prism. Internal PDF.
\bibitem{MO} Moonshine Operator. Internal PDF.
\bibitem{SPASC} SPASC. Internal PDF.
\bibitem{LCA} Low-Complexity Attractor. Internal PDF.
\bibitem{PIRTM} PIRTM. Internal PDF.
\bibitem{AutomorphicLearning} Automorphic Learning. Internal PDF.
\end{thebibliography}

\end{document}
